%%%%%%%%%%%%%%%%%
% Akhil Xavier — Resume
% Built with AltaCV
%%%%%%%%%%%%%%%%%

\documentclass[10pt,a4paper,ragged2e,withhyper]{altacv}

% -------------------------------------------------------------------
% Page layout
% -------------------------------------------------------------------

\geometry{
  left=1.25cm,
  right=1.25cm,
  top=1.5cm,
  bottom=1.5cm,
  columnsep=1.2cm
}

% Parallel columns
\usepackage{paracol}

% Header/footer support
\usepackage{fancyhdr}

% -------------------------------------------------------------------
% Fonts
% -------------------------------------------------------------------

\ifxetexorluatex
  \setmainfont{Lato}
\else
  \usepackage[default]{lato}
\fi

% -------------------------------------------------------------------
% Colors
% -------------------------------------------------------------------

\definecolor{VividPurple}{HTML}{3E0097}
\definecolor{SlateGrey}{HTML}{2E2E2E}
\definecolor{LightGrey}{HTML}{666666}

\colorlet{tagline}{VividPurple}
\colorlet{heading}{VividPurple}
\colorlet{headingrule}{VividPurple}
\colorlet{accent}{VividPurple}
\colorlet{emphasis}{SlateGrey}
\colorlet{body}{LightGrey}

% -------------------------------------------------------------------
% CV formatting
% -------------------------------------------------------------------

\renewcommand{\itemmarker}{{\small\textbullet}}
\renewcommand{\ratingmarker}{\faCircle}

% -------------------------------------------------------------------
% Document
% -------------------------------------------------------------------

\begin{document}

% -------------------------------------------------------------------
% Header
% -------------------------------------------------------------------

\name{Akhil Xavier}
\tagline{Embedded Software Engineer}

\photoR{2.5cm}{AkhilXavier.jpg}

\personalinfo{
  \email{akhilxavier@hotmail.com}
  \phone{+1 720 883 9760}
  \location{Denver, CO, USA}
  \linkedin{akhilxavier}
}

\makecvheader

% -------------------------------------------------------------------
% Layout
% -------------------------------------------------------------------

\AtBeginEnvironment{itemize}{\small}

% Equal-width columns
\columnratio{0.5}

% Start two-column layout
\begin{paracol}{2}

\cvsection{Work Experience}

\cvevent{Software Engineer, Senior}{Acuity Intelligent Spaces (post acquisition)}{Jan 2025- Ongoing}{Boulder, CO}
\begin{itemize}
    \item Developed LT6711GXE (HDMI to DP) bridge driver on Intel ADL platform (6.6 kernel).
    \item Development and customization on i915 DRM sub-system.
    \item Multiple development, validation, test activities on new QSYS NVM device.
\end{itemize}

\divider

\cvevent{Embedded Software Engineer}{QSC LLC (acquired by Acuity Intelligent Spaces)}{May 2023 - Dec 2024 }{Boulder, CO}
\begin{itemize}
    \item Enabled hardware accelerated QSV codec stack on Intel ADL platform.
    \item Optimized and validated hardware accelerated encode decode performance, GPU performance, video post processing (VPP) performance on intel ADL platform.
    \item Bring up of new encoder-decoder hardware units for a video streaming project.
\end{itemize}

\divider

\cvevent{Software Engineer}{Mersive Technologies Inc}{Nov 2019 – March 2023 }{Denver, CO}
\hspace{0.25cm}\textbf{(Nov 2019 - Dec 2022 India)}
\begin{itemize}
    \item Developed end-to-end HDMI input stack (TMDS-to-MIPI CSI), including driver, framework uEvent listener, and preview application.
    \item Implemented MI2S audio on Qualcomm SoC platforms.
    \item Improved CPU performance and thermal efficiency.
    \item Brought up multiple Mersive Pod hardware platforms.
\end{itemize}

\divider

\cvevent{Senior Software Engineer}{Inforce Computing India Pvt Ltd}{July 2015 – Oct 2019}{Bangalore, India}
\begin{itemize}
    \item Developed custom drivers for Bayer/YUV cameras, displays, chargers, HDMI-to-MIPI CSI, and DSI-to-HDMI ICs.
    \item Customized Qualcomm camera pipelines.
    \item QCOM SoC-based SBC/SOM bring-up and BSP development.
\end{itemize}

\divider

\cvevent{Software Intern Engineer}{Inforce Computing India Pvt Ltd}{Nov 2014 – July 2015}{Bangalore, India}
\begin{itemize}
    \item Qualcomm SOC based SBC,SOM bringup and BSP.
\end{itemize}

\switchcolumn

\cvsection{SKILLS AND EXPERTISE}
\begin{itemize}
    \item Languages: C, C++, Java (Intermediate Level)
    \item Android OS, Linux distributions (Linaro), Yocto Project
    \item GStreamer, ffmpeg
    \item HW Standards: MIPI CSI-2, MIPI-DSI, PCIe, HDMI, I2S, USB, CCI, DDC, I2C, SPI, UART
    \item Analyzing schematics and datasheets
    \item Chipsets: APQ8064, APQ8084, APQ8096, APQ8016, SDM660, Intel ADL
    \item Modules: Camera, Audio, Display, Charger, LowSpeed Peripherals, Image Processing
    \item BSP development, Linux Device Drivers
\end{itemize}

\cvsection{Education}

\cvevent{M Tech, Sensor System Technology}{VIT University, Vellore}{2013 -- 2015}{Tamil Nadu}

\divider

\cvevent{B.Tech, Electronics \& Communication Engineering}{Cochin University of Science and Technology, Kerala}{2007 -- 2011}{Cochin}

\cvsection{ACHIEVEMENTS}
\begin{itemize}
    \item Received best paper award during International Science and Technology Conference 2014 held at VIT University.
    \item Best employee of the month, Mersive Technologies.
    \item Elected as College Senate member.
    \item The Project submitted for ESNC 2014 conducted by European Union Space Agency was shortlisted to final stage.
    \item IEEE paper in EESCO “ISBN:978-1-4799-7676-8” on People Monitoring system.
\end{itemize}

\cvsection{TASKS ACCOMPLISHED}

\cvsubsection{Driver development and bring-up on QSC hardware}
\begin{itemize}
    \item Platform’s : QSC NVM
    \item SOC : Intel ADL N
    \item OS’s : QSYS OS, Linux kernel 6.6, Linux Kernel 5.15
    \item LT6711GXE HDMI to DP bridge I2C driver developed on QSC NVM device for dual input channels. The DP output flow to \\
          a video processing FPGA and then to SOC through PCIe.
    \item Enabled and validated hardware accelerated encoding, decoding, video post processing. GPU performance validated. \\
          Various performance metrics collected.

\newpage
\switchcolumn

    \item Encoding decoding performance on multiple pixel formats P010LE, YUV422, NV12, YUV444.
    \item Metrics on thermal performance and power consumption done.
    \item Enabled device mode operation on USB-C.
    \item Enabled native out HDMI HDCP support on Linux Intel ADL through libdrm.
    \item Bring up of various low speed peripherals done.
\end{itemize}

\cvsubsection{Multiple porting and customization of Mersive Pods}
\begin{itemize}
    \item Platform’s : MER6641
    \item SOC : APQ8084, APQ8096
    \item OS’s : Lollipop, Nougat, Oreo, Pi, Debian Linaro, Yocto build
    \item The Mersive pod is ported into upstream Android releases which are supported by the SOC Manufacturer. Also \\
          ported the Android based Mersive Pod into Debian based Linaro distribution.
    \item Customized pods peripherals based on customer usecase.
\end{itemize}

\cvsubsection{Driver development and Bring-up of various camera modules across multiple Android and Debian platforms.}
\begin{itemize}
    \item Platform’s : IFC6540, IFC6501, IFC6640, IFC6601, IFC6560
    \item SOC : APQ8084, APQ8096, APQ8016, SDM660, SDA845
    \item Sensor Module : OV5640 (YUV 5MP), OV7251, (Bayer Monochrome VGA), OV5645 (YUV 5MP), \\
                          IMX135 (Bayer 13MP), IMX230 (Bayer 21MP), OV8856 (Bayer 8MP), IMX290, OV8856, IMX258.
    \item OS’s : Lollipop, Marsh-mellow, Nougat, kernel 3.10 (Debian), Kernel 4.14 (Debian), Oreo PI, kernel 4.4
    \item Driver development and bring up of different imaging sensors are done. Different YUV, Bayer, \\
          Monochrome camera sensors are brought up.
    \item \href{www.github.com/akhilxavi/debian4.14-imx230driver}{www.github.com/akhilxavi/debian4.14-imx230driver}
\end{itemize}

\cvsubsection{Driver development and bring-up of HDMI to CSI convertor IC TC358840 and ADV7782}
\begin{itemize}
    \item Platform’s : IFC6540, IFC6501, IFC6640, IFC6601+, Mersive Pod
    \item SOC \& FPGA : APQ8084, APQ8096, TC358840 (4K HDMI to MIPI CSI-2 convertor), ACC1s80
    \item Camera Layers : Kernel, server process, media-controller, Android HAL, V4L2
    \item OS’s : Lollipop, Nougat
    \item Implemented HDMI-to-MIPI CSI via TC358840/ADV7782 FPGA, with uEvent-based hot-plug, resolution/FPS detection \\
          and dynamic camera pipeline reconfiguration.
\end{itemize}

\cvsubsection{Bring up, Feature implementation, BSP creation and bug Fixing across multiple projects and other miscellaneous works.}
\begin{itemize}
    \item Platform’s : IFC6540, IFC6501, IFC6640, IFC6601, IFC6309
    \item SOC \& FPGA : APQ8084, APQ8096, APQ8016, TC358840
    \item ACC and HW : ACC-1S80, Si5347
    \item OS’s : Lollipop, Nougat, Marshmallow, KitKat, Debian 3.10 kernel, Debian 4.14, Oreo

\switchcolumn

    \item Bug Fix for High Bit-rate video causing system shutdown.
    \item Fixed the issue of FD (file descriptor) leak in camera pipeline during resolution switch.
    \item Implemented device-tree overlay mechanism on 4.14 Debian kernel for using multiple sensors on same CSI port.
    \item \href{www.github.com/akhilxavi/debian4.14-DTOverlay}{www.github.com/akhilxavi/debian4.14-DTOverlay}
    \item  Bug fix for memory leak in camera pipeline.
    \item  Developed CCI/I2C read/write command line utilities to peek-poke camera sensor registers\\
           Communicate via IOCTL.
    \item \href{www.github.com/akhilxavi/CCI-I2C-R-W-utilities}{www.github.com/akhilxavi/CCI-I2C-R-W-utilities}
    \item  Porting of IFC6410 from KitKat to lollipop.
    \item  Xilinx DS890 ASIC based MIPI CSI2 frames generator bring up.
    \item USB Cameras.
\end{itemize}

\cvsubsection{Customization of camera Image Processing pipeline across various Platforms/Projects}
\begin{itemize}
    \item Platform’s : IFC6540, IFC6501, IFC6640, IFC6601
    \item SOC's : APQ8084, APQ8096, APQ8016
    \item ACC Card : ACC-1S80
    \item OS’s : Lollipop, Nougat
    \item Customized and debugged camera HAL across platforms, enabling dual/simultaneous camera, high-FPS, Y8, \\
          YUV422→YUV420, Bayer raw dump, CSI 4K frame sync, and stitched-frame pipeline; resolved memory and FD leaks.
\end{itemize}

\cvsubsection{Bring up of Analog and Digital audio across multiple platforms}
\begin{itemize}
    \item Platform’s : IFC6540, IFC6501, IFC6640, IFC6601
    \item SOC : APQ8084, APQ8096, APQ8016, TC358840
    \item ACC : ACC-1S80, WCD9330(codec), WCD9320(codec), PCM5100a(DAC)
    \item OS’s : Lollipop, Nougat, Marshmallow
    \item Brought up analog WCD codec audio for mic/HDMI stereo recording (44.1/48 kHz) and digital MI2S audio via SoC aDSP. \\
          \href{www.github.com/akhilxavi/android-3.18-MI2S-audio}{www.github.com/akhilxavi/android-3.18-MI2S-audio}
\end{itemize}

\cvsubsection{Multiple tasks across various platforms and projects}
\begin{itemize}
    \item Brought up Pico projection device BML050; enabled non-standard HDMI resolutions.
    (1920×1200, 800×600, 854×480) using custom EDID and PLL configuration.
    \item Configured GPIOs for PWM, I/O, and expansion-connector use cases.
    \item Configured PMIC regulators and LDOs for multiple peripherals.
    \item Enabled portrait-mode full-screen camera preview on TV displays.
    \item Brought up and verified Electronic Image Stabilization (EIS) using SSC.
    \item Reviewed release notes, test guides, and factory-support documentation during production.
    \item Improved Mersive Pod performance and implemented external USB caching.
    \item Resolved multiple software/hardware issues on Mersive Pod.
    \item Validated hardware platforms for next-generation Mersive products.
\end{itemize}

\end{paracol}
\end{document}
